\section{Performance Beyond Proportional Compute}
\label{sec:evidence}

The promise of allocation intelligence must be visible from outside the
system. Fusion should not require a customer to choose between frontier-level
capability and practical deployment efficiency. It should move the frontier
itself.

\subsection{A performance--compute frontier across three capability regimes}

Our current evaluation spans scientific reasoning, policy-constrained tool use,
and agentic coding through GPQA Diamond, $\tau^3$-Banking, and Terminal-Bench
v2.1, respectively \citep{rein2024gpqa,shi2026tauknowledge,merrill2026terminalbench}.
These workloads expose different failure modes, but ask the same systems
question: how much capability can be delivered for the compute consumed?
Table~\ref{tab:pareto} reports both score and the cost in US dollars of
completing each full benchmark run.

\begin{table}[H]
\centering
\caption{Reported score and full benchmark-run cost. Lower cost and higher
score are better. Within each comparison cohort, the highest score and lowest
cost are bold. Per-run costs are rounded to the nearest dollar.}
\label{tab:pareto}
\scriptsize
\begin{tabular}{@{}l
  r@{\hspace{1.2em}}r
  @{\hspace{2.0em}}r@{\hspace{1.2em}}r
  @{\hspace{2.0em}}r@{\hspace{1.2em}}r@{}}
\toprule
& \multicolumn{2}{c}{GPQA Diamond}
& \multicolumn{2}{c}{$\tau^3$-Banking}
& \multicolumn{2}{c}{Terminal-Bench v2.1} \\
\cmidrule(lr){2-3}\cmidrule(lr){4-5}\cmidrule(lr){6-7}
System & Score & Cost & Score & Cost & Score & Cost \\
\midrule
\textbf{Fusion Max} & \textbf{94.9\%} & \textbf{13}
  & \textbf{49.5\%} & \textbf{145} & 86.5\% & \textbf{74} \\
GPT 5.6 Sol (max) & 94.1\% & 22 & 44.3\% & 251
  & \textbf{88.0\%} & 86 \\
Claude Fable 5 (fallback) & 92.6\% & 45 & 38.1\% & 322
  & 84.6\% & 299 \\
Sakana Fugu Ultra & 94.4\% & 95 & 41.2\% & 224
  & 79.8\% & 144 \\
\addlinespace
\textbf{Fusion Code} & \textbf{93.5\%} & \textbf{7}
  & \textbf{48.5\%} & 159 & \textbf{89.9\%} & 71 \\
Claude Opus 5 (max) & 93.2\% & 18 & 42.1\% & 193
  & 89.1\% & 215 \\
Kimi K3 & \textbf{93.5\%} & 30 & 46.0\% & \textbf{122}
  & 85.0\% & \textbf{56} \\
\bottomrule
\end{tabular}
\end{table}

Figure~\ref{fig:observed-pareto} disaggregates the same results into native
benchmark coordinates. Each panel plots reported score against full
benchmark-run cost; no benchmark weighting or cross-benchmark normalization is
introduced. Each panel uses a linear cost axis fitted to that benchmark's
observed range, so equal horizontal distances represent equal cost differences
within the panel. In every panel, the preferred direction is toward the upper
left.

\begin{figure}[H]
\centering
\begin{tikzpicture}[x=1cm,y=1cm,font=\sffamily\scriptsize]
  \tikzset{
    pointlabel/.style={
      font=\sffamily\fontsize{6.0}{6.5}\selectfont,
      text=fusiongray,
      inner sep=0.6pt
    },
    fusionlabel/.style={
      font=\sffamily\bfseries\fontsize{6.0}{6.5}\selectfont,
      text=fusionink,
      inner sep=0.6pt
    }
  }
  % Shared legend and direction of improvement.
  \fill[fusionblue] (1.72,15.70) circle (2.4pt);
  \node[anchor=west,text=fusionink] at (1.90,15.70) {Fusion models};
  \filldraw[fill=white,draw=fusiongray,line width=0.7pt]
    (5.00,15.70) circle (2.2pt);
  \node[anchor=west,text=fusiongray] at (5.18,15.70) {Other models};
  \draw[fusionblue,densely dotted,line width=1.2pt]
    (8.05,15.70) -- (8.70,15.70);
  \node[anchor=west,text=fusionink] at (8.85,15.70)
    {Observed Pareto frontier};

  % Each benchmark has a complete linear cost axis.
  \foreach \xx/\label in {1.55/0,4.10/20,6.65/40,9.20/60,11.75/80,14.30/100} {
    \draw[fusiongray!14] (\xx,11.45) -- (\xx,14.65);
    \draw[fusiongray!45] (\xx,11.40) -- (\xx,11.50);
    \node[anchor=north,text=fusiongray] at (\xx,11.36) {\$\label};
  }
  \foreach \xx/\label in {1.55/100,4.10/150,6.65/200,9.20/250,11.75/300,14.30/350} {
    \draw[fusiongray!14] (\xx,6.95) -- (\xx,10.15);
    \draw[fusiongray!45] (\xx,6.90) -- (\xx,7.00);
    \node[anchor=north,text=fusiongray] at (\xx,6.86) {\$\label};
  }
  \foreach \xx/\label in {1.55/50,3.675/100,5.80/150,7.925/200,10.05/250,12.175/300,14.30/350} {
    \draw[fusiongray!14] (\xx,2.45) -- (\xx,5.65);
    \draw[fusiongray!45] (\xx,2.40) -- (\xx,2.50);
    \node[anchor=north,text=fusiongray] at (\xx,2.36) {\$\label};
  }

  % Panel frames, horizontal grids, and native score ticks.
  \draw[fusiongray!35] (1.55,11.45) rectangle (14.30,14.65);
  \foreach \yy/\score in {11.450/91.5,11.850/92.0,12.250/92.5,12.650/93.0,13.050/93.5,13.450/94.0,13.850/94.5,14.250/95.0,14.650/95.5} {
    \draw[fusiongray!18] (1.55,\yy) -- (14.30,\yy);
    \node[anchor=east,text=fusiongray] at (1.42,\yy) {\score\%};
  }
  \node[anchor=west,text=fusionink] at (1.55,14.92)
    {\textbf{(a) GPQA Diamond}};

  \draw[fusiongray!35] (1.55,6.95) rectangle (14.30,10.15);
  \foreach \yy/\score in {6.950/35.0,7.407/37.5,7.864/40.0,8.321/42.5,8.779/45.0,9.236/47.5,9.693/50.0,10.150/52.5} {
    \draw[fusiongray!18] (1.55,\yy) -- (14.30,\yy);
    \node[anchor=east,text=fusiongray] at (1.42,\yy) {\score\%};
  }
  \node[anchor=west,text=fusionink] at (1.55,10.42)
    {\textbf{(b) $\tau^3$-Banking}};

  \draw[fusiongray!35] (1.55,2.45) rectangle (14.30,5.65);
  \foreach \yy/\score in {2.450/78,2.907/80,3.364/82,3.821/84,4.279/86,4.736/88,5.193/90,5.650/92} {
    \draw[fusiongray!18] (1.55,\yy) -- (14.30,\yy);
    \node[anchor=east,text=fusiongray] at (1.42,\yy) {\score\%};
  }
  \node[anchor=west,text=fusionink] at (1.55,5.92)
    {\textbf{(c) Terminal-Bench v2.1}};

  % GPQA Diamond frontiers.
  \draw[fusionblue,densely dotted,line width=1.2pt,line cap=round]
    (2.443,13.050) -- (3.208,14.170);

  \filldraw[fill=white,draw=fusiongray,line width=0.7pt]
    (4.355,13.530) circle (2.1pt);
  \filldraw[fill=white,draw=fusiongray,line width=0.7pt]
    (7.288,12.330) circle (2.1pt);
  \filldraw[fill=white,draw=fusiongray,line width=0.7pt]
    (13.663,13.770) circle (2.1pt);
  \filldraw[fill=white,draw=fusiongray,line width=0.7pt]
    (3.845,12.810) circle (2.1pt);
  \filldraw[fill=white,draw=fusiongray,line width=0.7pt]
    (5.375,13.050) circle (2.1pt);
  \fill[fusionblue] (3.208,14.170) circle (2.5pt);
  \node[diamond,draw=fusionblue,fill=fusionblue,inner sep=1.45pt]
    at (2.443,13.050) {};

  \node[fusionlabel,anchor=east] at (2.38,12.99) {Fusion Code};
  \node[fusionlabel,anchor=west] at (3.33,14.17) {Fusion Max};
  \node[pointlabel,anchor=east] at (3.72,12.81) {Claude Opus 5};
  \node[pointlabel,anchor=west] at (4.48,13.53) {GPT 5.6 Sol};
  \node[pointlabel,anchor=west] at (5.50,13.05) {Kimi K3};
  \node[pointlabel,anchor=west] at (7.41,12.33) {Claude Fable 5};
  \node[pointlabel,anchor=east] at (13.54,13.77) {Sakana Fugu Ultra};

  % Tau3-Banking frontiers.
  \draw[fusionblue,densely dotted,line width=1.2pt,line cap=round]
    (2.672,8.961) -- (3.845,9.601);

  \filldraw[fill=white,draw=fusiongray,line width=0.7pt]
    (9.251,8.651) circle (2.1pt);
  \filldraw[fill=white,draw=fusiongray,line width=0.7pt]
    (12.872,7.517) circle (2.1pt);
  \filldraw[fill=white,draw=fusiongray,line width=0.7pt]
    (7.874,8.084) circle (2.1pt);
  \filldraw[fill=white,draw=fusiongray,line width=0.7pt]
    (6.293,8.248) circle (2.1pt);
  \filldraw[fill=white,draw=fusiongray,line width=0.7pt]
    (2.672,8.961) circle (2.1pt);
  \fill[fusionblue] (3.845,9.601) circle (2.5pt);
  \node[diamond,draw=fusionblue,fill=fusionblue,inner sep=1.45pt]
    at (4.559,9.419) {};

  \node[pointlabel,anchor=east] at (2.55,8.96) {Kimi K3};
  \node[fusionlabel,anchor=south east] at (3.73,9.73) {Fusion Max};
  \node[fusionlabel,anchor=west] at (4.69,9.42) {Fusion Code};
  \node[pointlabel,anchor=east] at (6.17,8.25) {Claude Opus 5};
  \node[pointlabel,anchor=west] at (7.99,8.08) {Sakana Fugu Ultra};
  \node[pointlabel,anchor=west] at (9.37,8.65) {GPT 5.6 Sol};
  \node[pointlabel,anchor=east] at (12.75,7.52) {Claude Fable 5};

  % Terminal-Bench frontiers.
  \draw[fusionblue,densely dotted,line width=1.2pt,line cap=round]
    (1.805,4.050) -- (2.443,5.170);

  \filldraw[fill=white,draw=fusiongray,line width=0.7pt]
    (3.080,4.736) circle (2.1pt);
  \filldraw[fill=white,draw=fusiongray,line width=0.7pt]
    (12.133,3.959) circle (2.1pt);
  \filldraw[fill=white,draw=fusiongray,line width=0.7pt]
    (5.545,2.861) circle (2.1pt);
  \filldraw[fill=white,draw=fusiongray,line width=0.7pt]
    (8.563,4.987) circle (2.1pt);
  \filldraw[fill=white,draw=fusiongray,line width=0.7pt]
    (1.805,4.050) circle (2.1pt);
  \fill[fusionblue] (2.570,4.393) circle (2.5pt);
  \node[diamond,draw=fusionblue,fill=fusionblue,inner sep=1.45pt]
    at (2.443,5.170) {};

  \node[pointlabel,anchor=west] at (1.92,3.94) {Kimi K3};
  \node[fusionlabel,anchor=west] at (2.56,5.25) {Fusion Code};
  \node[fusionlabel,anchor=west] at (2.69,4.31) {Fusion Max};
  \node[pointlabel,anchor=west] at (3.20,4.74) {GPT 5.6 Sol};
  \node[pointlabel,anchor=west] at (5.66,2.86) {Sakana Fugu Ultra};
  \node[pointlabel,anchor=west] at (8.68,4.99) {Claude Opus 5};
  \node[pointlabel,anchor=west] at (12.25,3.96) {Claude Fable 5};

  \node[rotate=90,anchor=south,text=fusiongray] at (0.23,8.55)
    {Native benchmark score};
  \node[anchor=north,text=fusiongray] at (7.93,1.65)
    {Full benchmark-run cost (USD)};
\end{tikzpicture}
\caption{Native score--cost frontiers for the three reported benchmarks. Each
point represents a full benchmark run from Table~\ref{tab:pareto}. Every panel
uses its native score and an independently scaled linear cost axis, so
positions should be interpreted within, not across, panels. Dotted blue lines
connect the observed non-dominated operating points after including Fusion;
they are not fitted curves. Exact values are reported in Table~\ref{tab:pareto}.}
\label{fig:observed-pareto}
\end{figure}

Fusion Max reduces the aggregate cost of this fixed three-benchmark portfolio
from \$359 to \$232 relative to GPT 5.6 Sol, a 35.4\% reduction. It scores
higher on GPQA Diamond and $\tau^3$-Banking and remains within 1.5 percentage
points on Terminal-Bench. Against Claude Fable 5, Fusion Max scores higher on
all three workloads while reducing each full-run cost by 55--75\%; aggregate
cost falls from \$666 to \$232.

Fusion Code presents a second operating point. It meets or exceeds the reported
score of Claude Opus 5 on all three workloads while reducing full-run cost by
18--67\%; aggregate cost falls from \$426 to \$237. Kimi K3 illustrates a
different trade-off: Fusion Code matches or improves its score on all three
workloads, but spends more on the Banking and Terminal-Bench runs. We therefore
present that comparison as a capability gain, not as a universal cost
reduction.

The result is a Pareto claim, not a discount claim. Fusion retains stronger
judgment where it is valuable and removes the assumption that every
intermediate unit of reasoning must be purchased at the same capability tier.
A system that is only cheaper because it silently accepts worse outcomes has
not demonstrated allocation intelligence.

\subsection{Why asymmetry changes the economics of capability}

The cost profile follows directly from Fusion's architecture. A frontier-heavy
orchestration design assembles its worker pool from several of the most capable
general-purpose models. This can increase quality, but it also duplicates the
most expensive class of computation across planning, generation, critique, and
verification. Fusion instead composes a wider capability spectrum. It asks
where frontier intelligence is decisive and where a specialized or efficient
source can add more information per unit of compute.

Sakana Fugu Ultra provides a useful external operating point
\citep{fugu2026}. It uses learned orchestration across a frontier-heavy worker
pool, while Fusion uses task-state understanding to allocate across an
asymmetric supply of models. In our available full-suite runs, it reaches
94.4\% on GPQA Diamond, 41.2\% on
$\tau^3$-Banking, and 79.8\% on Terminal-Bench, at full-run costs of \$94.55,
\$224.14, and \$143.61. Its GPQA result is 185/196 scored cases; two additional
selected cases ended in execution errors.

These results do not support a blanket order-of-magnitude cost claim. Across
the fixed three-benchmark portfolio, Sakana Fugu Ultra costs \$462.30, compared
with \$232 for Fusion Max: Fusion Max costs about 49.8\% less while scoring
higher on all three workloads. Costs differ by roughly an order of magnitude
on one workload, not across the portfolio. On GPQA, where the two systems are
within 0.5 percentage points, Sakana Fugu Ultra costs 7.3 times as much as
Fusion Max and 13.5 times as much as Fusion Code. The Banking and
Terminal-Bench runs instead show a joint capability and cost advantage for
Fusion, rather than cost parity at matched quality.

The comparison is directional rather than perfectly controlled: model and
provider snapshots differ, the Banking and Terminal-Bench runs are
official-compatible rather than identical submissions, and the Terminal-Bench
run is not leaderboard-comparable. Even within that boundary, the architectural
contrast remains useful. One system primarily orchestrates frontier
intelligence; the other learns where frontier intelligence is necessary.

The point is not that an asymmetric pool is inherently superior. A weak source
that adds misleading or redundant evidence is not efficient at any price
\citep{li2025rethinkingmoa}. The advantage appears only when the system has the
capability understanding and task-state awareness to use asymmetry
selectively. That allocation layer---not the mere presence of smaller
models---is the source of the frontier shift.

\subsection{Effectiveness that survives the workflow}

One-shot benchmarks capture only part of Fusion's value. In an agentic system,
the demand for intelligence changes after every observation and action. A
design that improves a single answer can still fail over a trajectory by
repeating work, losing state, acting on incompatible plans, or declaring
success before verification.

Fusion is designed so that its benefit persists across the workflow. The same
canonical state, uncertainty model, and authority boundary apply whether the
system is investigating, editing, testing, recovering from failure, or
committing a result. The relevant outcome becomes successful task completion,
not the quality of an isolated response. The relevant cost becomes compute per
successful trajectory, not price per API call.

\begin{tcolorbox}[
  enhanced,
  breakable,
  colback=fusionpaper,
  colframe=fusionpaper,
  boxrule=0pt,
  arc=2pt,
  left=11pt,right=11pt,top=9pt,bottom=9pt
]
{\sffamily\bfseries\color{fusionink}A LONG-HORIZON TRACE: ADAPTATION,
NOT REPETITION}\par\smallskip
In the \emph{path-tracing} task from Fusion Code's full Terminal-Bench v2.1
run, the system had to reconstruct an image produced by an unknown binary as a
dependency-free C program whose compressed source had to satisfy a strict size
limit. Its initial inspection path failed because basic utilities were absent.
It changed methods, using available scripting and binary-analysis tools to
inspect the image, analyze the binary, recover the underlying scene, implement
a compact renderer, compile it, and compare the generated output against the
target.

The trajectory ran for about 29 minutes across 49 agent steps. Before
completion,
the system measured a normalized L2 similarity of 0.999998; the benchmark's
independent verifier then passed all five tests, including compilation,
dependency, execution, and similarity checks. This single trace is not an
aggregate claim. It illustrates a property that aggregate evaluation must test:
when execution invalidates an assumption, the system can preserve state,
reduce the new uncertainty, change its next action, and still close the loop
with external verification.
\end{tcolorbox}

We therefore evaluate long-horizon behavior in terms of end-to-end completion,
total cost per success, frontier calls per trajectory, state consistency,
recovery after failed actions, premature completion, and degradation as
workflow length increases. This is where allocation intelligence becomes a
system capability rather than a routing optimization.

\subsection{A transparent standard of evidence}

Public comparisons should expose the full operating point: task set, model and
policy versions, scorer, attempts, timeout, concurrency, cache treatment,
latency, and all internal model usage. Client-visible tokens are not the total
cost of orchestration. Infrastructure and scorer failures must remain separate
from system failures. Fumo Lab will report these boundaries explicitly,
especially where a fast-moving external system cannot be reproduced under a
perfectly matched protocol.
